%% kapa: exam=1 round=20 year=2009 period=2 form=A subject=appraisal_law subject_ko=부동산관계법규 seq=25 number=65
%% source: Q-Net 감정평가사 기출문제 (cst00309) artlSeq=5024178 file=extracted/1008/5024178/2009년 제20회 감정평가사 1차시험 기출문제/2교시.hwp
%% conversion: kordoc 4.24.51 -> tools/gen1.py (auto); review=proofread-sonnet; flags=none
\begin{question}{65}
국유재산법령상 국유재산에 관한 설명으로 옳지 \uline{않은} 것은?
\begin{choices}
  \item 국가의 부담에 의하여 국유로 된 부동산과 그 종물은 국유재산이다.
  \item 기부의 채납에 의하여 국유로 된 선박․항공기와 그들의 종물은 국유재산이다.
  \item 법령이나 조약의 규정에 의하여 국유로 된 부동산신탁의 수익권은 국유재산이다.
  \item 국유재산은 그 종류와 용도에 무관하게 공법의 적용을 받는 공물에 해당한다.
  \item 국가가 직접 그 사무용ㆍ사업용 또는 공무원의 주거용으로 사용하거나 사용하기로 결정한 재산은 공용재산이다.
\end{choices}
\end{question}
\answer{4}
